\clearpage
\renewcommand{\sectioncolor}{slate}
\section{Appendix: dated source ledger}
\markboth{Source ledger}{}

Every substantive claim in this dossier, with its origin. Contested items are marked.

\begin{center}\small
\setlength{\tabcolsep}{5pt}
\begin{tabular}{@{}>{\raggedright\arraybackslash}p{1.35cm}>{\raggedright\arraybackslash}p{6.0cm}>{\raggedright\arraybackslash}p{9.2cm}@{}}
\toprule
\rowcolor{slate!14}
\textbf{Year} & \textbf{Result} & \textbf{Who / where, and status}\\
\midrule
\multicolumn{3}{@{}l}{\textbf{\color{quantum}Escapes past intuition}}\\
1829--54 & Non-Euclidean geometry & Lobachevsky 1829; Bolyai 1832; Riemann 1854\\
1868 & Pseudosphere model (the re-grounding) & Beltrami; later Klein and Poincaré disk models\\
1932 & Hilbert-space formulation of QM & von Neumann, \textit{Math.\ Grundlagen der Quantenmechanik}\\
1936 & Non-distributive quantum logic & Birkhoff \& von Neumann, \textit{Annals of Mathematics}\\
1980s--94 & Non-commutative geometry & Connes; \textit{Noncommutative Geometry}, 1994\\
2011 & Contextuality as a sheaf-theoretic obstruction & Abramsky \& Brandenburger, \textit{New J.\ Phys.}\\
\midrule
\multicolumn{3}{@{}l}{\textbf{\color{bio}Diagrammatic re-grounding (Move 3)}}\\
1948 & Feynman diagrams & Feynman\\
1971 & Graphical tensor notation & Penrose\\
2008--11 & ZX-calculus & Coecke \& Duncan; complete for qubit QM\\
\midrule
\multicolumn{3}{@{}l}{\textbf{\color{sand}Instruments that generated mathematics (Move 4)}}\\
1981 & Scanning tunnelling microscope & Binnig \& Rohrer; Nobel Prize 1986\\
1986 & Atomic force microscope & Binnig, Quate \& Gerber\\
1978 & Forced-unbinding model & Bell, \textit{Science}\\
1997 & Dynamic force spectroscopy theory & Evans \& Ritchie, \textit{Biophys.\ J.}\\
\midrule
\multicolumn{3}{@{}l}{\textbf{\color{nano}Mathematics that already works for molecular systems}}\\
1972-- & Deficiency-zero / deficiency-one theorems & Horn \& Jackson 1972; Feinberg onwards\\
1974--75 & Kinetic proofreading & Hopfield (PNAS) 1974; Ninio 1975\\
1976--77 & Stochastic simulation algorithm & Gillespie\\
1976 & Lindblad master equation & Lindblad; Gorini--Kossakowski--Sudarshan same year\\
1989 & Hierarchical equations of motion (HEOM) & Tanimura \& Kubo\\
1997 & Non-equilibrium work relation & Jarzynski, \textit{Phys.\ Rev.\ Lett.}\\
1999 & Fluctuation theorem & Crooks, \textit{Phys.\ Rev.\ E}\\
2002--09 & Persistent homology / TDA & Edelsbrunner, Letscher \& Zomorodian 2002; Carlsson 2009\\
2015 & Thermodynamic uncertainty relation & Barato \& Seifert, \textit{Phys.\ Rev.\ Lett.}\\
2017 & Compositional framework for reaction networks & Baez \& Pollard\\
\midrule
\multicolumn{3}{@{}l}{\textbf{\color{ember}Machine-extended mathematics (Move 5)}}\\
2019 & Boltzmann generators & Noé et al., \textit{Science}\\
2020 & Symbolic regression for physical law & Udrescu \& Tegmark (AI~Feynman)\\
2020--22 & Liquid Tensor Experiment & Scholze posed Dec 2020; completed in Lean Jul 2022\\
2021 & AI-guided mathematical intuition & Davies et al., \textit{Nature} --- knot theory and
 Kazhdan--Lusztig conjectures, both subsequently proved\\
2021 & AlphaFold2 & Jumper et al., \textit{Nature}\\
2021 & cryoDRGN, continuous conformational heterogeneity & Zhong et al., \textit{Nature Methods}\\
\midrule
\multicolumn{3}{@{}l}{\textbf{\color{crimson}Quantum biology --- graded}}\\
1978 & Radical-pair magnetoreception proposed & Schulten et al. \;\textsc{strong, not closed}\\
2000 & Decoherence argument against Orch-OR & Tegmark, \textit{Phys.\ Rev.\ E} \;\textsc{objection stands}\\
2007 & ``Long-lived coherence'' in light harvesting & Engel et al., \textit{Nature}
 \;\textsc{substantially revised}\\
2016 & Magnetoreception review & Hore \& Mouritsen, \textit{Annu.\ Rev.\ Biophys.}\\
2021 & Cryptochrome-4 magnetic sensitivity \emph{in vitro} & Xu et al., \textit{Nature}\\
\bottomrule
\end{tabular}
\end{center}

\subsection{Where to start tomorrow}

If you want one entry point into each of the three levers, these are the ones I would pick for a
reader with your background:

\begin{center}\small
\begin{tabular}{@{}>{\raggedright\arraybackslash}p{3.0cm}>{\raggedright\arraybackslash}p{13.6cm}@{}}
\toprule
\rowcolor{slate!14}
\textbf{Lever} & \textbf{Entry point}\\
\midrule
\textbf{\color{sand}1 · Body} & Single-molecule force spectroscopy: read Bell 1978, then Evans \&
 Ritchie 1997, then Jarzynski 1997. Three short papers, and you will have watched an instrument
 create a mathematics.\\
\textbf{\color{quantum}2 · Re-ground} & Coecke \& Kissinger's diagrammatic treatment of quantum
 theory. Start with the ZX rewrite rules and prove something small by pushing pictures.\\
\textbf{\color{ember}3 · Externalise} & Lean and \texttt{mathlib}. Formalise one theorem you already
 understand --- the point is to feel the difference between verified and understood.\\
\rowcolor{bio!8}
\textbf{\color{bio}The audit} & Feinberg's deficiency-zero theorem. It is the cleanest existing example
 of mathematics built to fit molecular systems rather than borrowed from physics.\\
\bottomrule
\end{tabular}
\end{center}

\vspace{12pt}
\begin{center}
\begin{tikzpicture}
\node[rounded corners=5pt,fill=bio!92!black,text=white,inner sep=13pt,text width=15.6cm,align=center]
 {\mbox{\large\bfseries Extend the body \ $\cdot$ \ Re-ground the abstraction \ $\cdot$ \ Externalise the bookkeeping}\\[7pt]
  \normalfont\small The embodiment of mathematics is not a ceiling. It is a specification of where
  the next lever has to be placed.};
\end{tikzpicture}
\end{center}
